\chapter{Studiu de piață și metode existente}

\section{Platforme educaționale existente}
Piața platformelor online de învățare a limbilor străine este destul de variată, platformele existente acoperind nevoi diferite ale utilizatorilor și funcționând pe modele de business diferite, fiecare adresând o parte din problemă fără a oferi o soluție completă pentru cineva care vrea atât materiale de studiu structurate, cât și posibilitatea de a lucra direct cu un profesor.

\textbf{Duolingo}

Duolingo este probabil cea mai cunoscută platformă de învățare a limbilor străine la nivel global, punând la dispoziția utilizatorilor lecții scurte bazate pe mecanisme de tip joc, organizate pe niveluri de dificultate și accesibile gratuit în varianta de bază, cu o versiune premium care adaugă funcționalități suplimentare și elimină reclamele, modelul platformei bazându-se pe încurajarea unui obicei zilnic de studiu prin recompense și streak-uri \cite{duolingo}. Această abordare o face potrivită pentru utilizatorii care vor să studieze câteva minute pe zi fără un angajament serios față de un program de studiu, platforma reușind să atragă un număr foarte mare de utilizatori tocmai prin accesibilitatea și simplitatea experienței pe care o oferă.

Limitele devin evidente atunci când utilizatorul dorește mai mult decât exerciții scurte și recompense virtuale, platforma neoferind niciun mecanism de interacțiune cu un profesor real, conținutul fiind standardizat și identic pentru toți utilizatorii fără a putea fi adaptat nevoilor individuale, iar sistemul de urmărire a progresului bazându-se pe indicatori gamificați care nu reflectă neapărat un nivel lingvistic real și funcțional.

\textbf{Coursera}

Coursera funcționează pe un model complet diferit față de Duolingo, punând accentul pe cursuri academice create de universități și instituții de prestigiu din întreaga lume, organizate în module cu materiale video, teme și evaluări, finalizarea unui curs ducând la obținerea unui certificat recunoscut care poate fi adăugat în CV-ul profesional al utilizatorului \cite{coursera}. Nivelul conținutului disponibil pe platformă este ridicat, Coursera adresându-se în principal utilizatorilor care vor să urmeze un curs structurat serios, nu neapărat să practice engleza conversațional sau să primească îndrumare personalizată din partea unui profesor.

Interacțiunea directă cu un profesor lipsește complet din modelul standard al platformei, comunicarea limitându-se la forumuri de discuții unde studenții interacționează între ei sau cu asistenți de curs, sesiunile live nefiind parte din oferta obișnuită, conținutul fiind preînregistrat și identic pentru toți utilizatorii indiferent de nivelul sau obiectivele lor specifice.

\textbf{iTalki}

iTalki funcționează ca o piață de meditații online, conectând studenți cu profesori sau vorbitori nativi disponibili pentru sesiuni individuale tarifate per sesiune, oferind o gamă largă de profesori pentru aproape orice limbă la prețuri foarte variate, ceea ce o face accesibilă financiar pentru mulți utilizatori \cite{italki}. Flexibilitatea este principalul avantaj al platformei, studenții putând alege profesori în funcție de preț, disponibilitate și recenziile primite de la alți utilizatori.

Problema principală a iTalki constă în faptul că nu oferă niciun plan de studiu propriu sau structurare a procesului de învățare, fiecare sesiune fiind practic independentă de celelalte, profesorul neavând acces la niciun fel de informație despre ce a studiat studentul până în acel moment și neputând adapta sesiunea în funcție de lacunele sau obiectivele specifice ale acestuia. În plus, lipsa unui proces formal de verificare a calificărilor profesorilor face ca oricine să poată crea un profil și oferi sesiuni, calitatea variind astfel considerabil de la un profesor la altul și lăsând studentul să evalueze singur dacă alegerea sa a fost bună.

\textbf{Preply}

Preply funcționează pe un model similar cu iTalki, conectând studenți cu profesori pentru sesiuni online individuale, dar adăugând față de aceasta un strat suplimentar de structurare, incluzând unele materiale de studiu integrate în platformă și un proces de selecție mai riguros pentru profesori, ceea ce o face o alegere mai sigură pentru utilizatorii care vor să lucreze cu cineva verificat \cite{preply}. Platforma urmărește și progresul studentului la nivel de sesiuni susținute, oferind o imagine mai clară a activității față de alte platforme de meditații online.

Chiar și cu aceste îmbunătățiri față de iTalki, materialele de studiu disponibile de Preply nu sunt integrate organic cu sesiunile live, rămânând mai degrabă resurse auxiliare decât un plan de studiu coerent urmărit sistematic împreună cu profesorul, iar progresul urmărit reflectând în principal numărul de sesiuni susținute, nu competențele efectiv dobândite prin parcurgerea unui conținut structurat. Din toate platformele analizate, Preply este cea mai apropiată de conceptul platformei propuse în această lucrare, dar distanța dintre ele rămâne semnificativă în ceea ce privește integrarea conținutului didactic cu sesiunile video.

\section{Analiză comparativă a platformelor}

Din analiza celor patru platforme se conturează clar o separare între cele două segmente identificate în capitolul 2, platformele de conținut neoferind interacțiune directă cu profesori, iar platformele de meditații online neavând un plan de studiu propriu integrat, niciuna dintre ele acoperind complet nevoile unui utilizator care dorește ambele lucruri din același loc.

\Cref{tab:platform-comparison} centralizează compararea celor patru platforme pe criteriile esențiale care asigură o experiență completă de învățare a limbii engleze, evidențiind golul pe care platforma propusă în această lucrare încearcă să îl acopere.

\begin{table}[htbp]
    \centering
    \resizebox{\textwidth}{!}{%
        \begin{tabular}{l c c c c}
            \toprule
            \textbf{Criteriu} & \textbf{Duolingo} & \textbf{Coursera} & \textbf{iTalki} & \textbf{Preply} \\
            \midrule
            Materiale de studiu structurate & Da & Da & Nu & Parțial \\
            Sesiuni video cu profesori & Nu & Nu & Da & Da \\
            Urmărire progres & Da & Da & Nu & Parțial \\
            Verificare profesori & Nu & Nu & Parțial & Da \\
            Model de plată & Gratuit/Premium & Abonament/Per curs & Per sesiune & Abonament \\
            \bottomrule
        \end{tabular}%
    }
    \caption{Analiza comparativă a platformelor educaționale existente}
    \label{tab:platform-comparison}
\end{table}

\section{Tehnologii utilizate și alternative}
Alegerea tehnologiilor pentru implementarea platformei a urmărit în principal potrivirea cu cerințele funcționale identificate și compatibilitatea componentelor între ele, fiecare decizie fiind luată după analizarea alternativelor disponibile și argumentată atât din punct de vedere calitativ, cât și cantitativ acolo unde a fost posibil.

\subsection{Frontend}
Pentru interfața de utilizator a fost ales Next.js, un framework bazat pe React, care îi extinde funcționalitățile de bază cu suport nativ pentru randare pe server, generare statică a paginilor și un sistem de rutare bazat pe structura de fișiere a proiectului, reducând semnificativ configurarea necesară față de o aplicație React obișnuită și oferind performanță mai bună la prima încărcare a paginilor \cite{nextjs}. Alegerea Next.js față de React a venit din nevoia de randare pe server a paginilor publice ale platformei, îmbunătățind experiența utilizatorului și comportamentul general al aplicației fără a sacrifica interactivitatea necesară secțiunilor dinamice.

Alternativa principală luată în considerare a fost Nuxt.js, meta-framework-ul construit peste Vue.js care oferă suport nativ pentru randare pe server, similar cu ce oferă Next.js pentru React, beneficiind de o curbă de învățare mai ușoară și o sintaxă considerată mai accesibilă față de React \cite{nuxtjs}. Decizia de a nu folosi Nuxt.js a venit din faptul că ecosistemul de biblioteci disponibil pentru React și Next.js este considerabil mai bogat, oferind mai multe opțiuni pentru funcționalitățile necesare platformei (\Cref{tab:frontend-comparison}).

\begin{table}[htbp]
    \centering
    \small
    \begin{tabular}{l c c}
        \toprule
        \textbf{Criteriu} & \textbf{Next.js} & \textbf{Nuxt.js} \\
        \midrule
        Randare pe server & Da & Da \\
        Limbaj principal & JavaScript/TypeScript & JavaScript/TypeScript \\
        Curbă de învățare & Medie & Mică \\
        Ecosistem biblioteci & Foarte mare & Mediu \\
        Rutare integrată & Da & Da \\
        \bottomrule
    \end{tabular}
    \caption{Compararea framework-urilor frontend}
    \label{tab:frontend-comparison}
\end{table}

\subsection{Backend}
Pentru server a fost ales Spring Boot, un framework bazat pe Java care simplifică configurarea aplicațiilor Spring, oferind funcționalități integrate pentru construirea de REST API-uri, gestionarea securității prin Spring Security și accesul la baza de date prin JPA, impunând totodată o structură clară a codului care face aplicația mai ușor de întreținut și de extins pe termen lung \cite{springboot}. Alegerea Spring Boot a venit în principal din nevoia de a gestiona funcționalități complexe ale aplicației, care implică roluri și permisiuni, fluxuri de aprobare, sisteme de credite și integrări cu servicii externe, toate acestea beneficiind de structura și instrumentele pe care le oferă ecosistemul Spring.

Principala alternativă considerată a fost Express.js, un framework minimalist pentru Node.js bazat pe JavaScript, apreciat pentru flexibilitatea sa și pentru rapiditatea cu care poate fi pornit un server funcțional, bazându-se pe un model asincron care îl face eficient pentru operații de tip I/O și potrivit pentru API-uri simple cu cerințe mai puțin complexe \cite{expressjs}. Lipsa unei structuri impuse poate fi un avantaj pentru proiecte mici, dar devine un dezavantaj pe măsură ce aplicația crește, codul riscând să devină greu de urmărit și de întreținut, iar absența unui sistem de securitate integrat comparabil cu Spring Security reprezentând un argument suplimentar în favoarea alegerii Spring Boot pentru nevoile acestei platforme (\Cref{tab:backend-comparison}).

\begin{table}[htbp]
    \centering
    \small
    \begin{tabular}{l c c}
        \toprule
        \textbf{Criteriu} & \textbf{Spring Boot} & \textbf{Express.js} \\
        \midrule
        Limbaj & Java & JavaScript/TypeScript \\
        Securitate integrată & Da & Nu \\
        Structură impusă & Da & Nu \\
        Model execuție & Multi-thread & Single-thread asincron \\
        Complexitate configurare & Medie & Mică \\
        \bottomrule
    \end{tabular}
    \caption{Compararea framework-urilor backend}
    \label{tab:backend-comparison}
\end{table}

\subsection{Bază de date}
Pentru stocarea datelor a fost aleasă PostgreSQL, o bază de date relațională open source cu suport pentru tranzacții complexe, integritate referențială și un limbaj SQL extins cu funcționalități avansate, alegerea unui model relațional venind firesc din structura datelor platformei, entitățile sistemului formând relații clare și bine definite între ele: utilizatorii, cursurile, lecțiile, sesiunile și tranzacțiile se leagă unele de altele în mod natural într-o schemă relațională \cite{postgresql}.

Alternativa analizată a fost MongoDB, o bază de date de tip NoSQL orientată pe documente cu o schemă flexibilă care permite câmpuri diferite de la un document la altul, funcționând bine pentru date nestructurate sau cu schemă variabilă și oferind o scalabilitate orizontală mai ușor de realizat față de bazele de date relaționale tradiționale \cite{mongodb}. Pentru platforma propusă, datele au o structură fixă și bine definită, iar operații precum deducerea creditelor la momentul rezervării necesită consistența și integritatea datelor pe care PostgreSQL le asigură nativ, MongoDB nefiind la fel de potrivit pentru acest tip de cerințe (\Cref{tab:database-comparison}).

\begin{table}[htbp]
    \centering
    \small
    \begin{tabular}{l c c}
        \toprule
        \textbf{Criteriu} & \textbf{PostgreSQL} & \textbf{MongoDB} \\
        \midrule
        Tip & Relațional & Non-relațional \\
        Schemă & Fixă & Flexibilă \\
        Relații complexe & Da & Limitat \\
        Consistență tranzacțională & Da & Parțial \\
        Open Source & Da & Da \\
        \bottomrule
    \end{tabular}
    \caption{Compararea soluțiilor de stocare a datelor}
    \label{tab:database-comparison}
\end{table}

\subsection{Integrări externe}
Pe lângă tehnologiile principale utilizate, platforma integrează trei servicii externe pentru funcționalități care ar fi necesitat un efort semnificativ de implementare de la zero și care sunt mai bine gestionate prin servicii specializate.

Sesiunile video sunt gestionate prin Daily.co, un serviciu care oferă o infrastructură completă de videoconferință accesibilă prin API, permițând crearea de camere de sesiune și desfășurarea apelurilor video direct în browser, fără ca utilizatorul să fie nevoit să instaleze o aplicație separată, întreaga interacțiune desfășurându-se în cadrul platformei \cite{dailyco}. Chiar dacă serviciile similare de videoconferință disponibile pe piață aveau funcționalități similare, Daily.co s-a diferențiat prin simplitatea integrării și prin faptul că oferă o interfață video completă care poate fi inclusă direct în aplicație, reducând semnificativ timpul de implementare.

Sistemul de plăți și abonamente este gestionat prin Stripe, un serviciu de procesare a plăților online cu suport pentru abonamente recurente, webhook-uri pentru notificarea serverului la evenimentele de plată și un portal dedicat prin care utilizatorii își pot gestiona singuri abonamentul fără a fi necesară intervenția unui administrator \cite{stripe}. Spre deosebire de celelalte soluții de procesare a plăților, Stripe oferă o documentație mai bine structurată, un API mai modern și un set complet de funcționalități pentru gestionarea abonamentelor, reprezentând astfel alegerea mai potrivită pentru nevoile platformei.

Migrările bazei de date sunt gestionate prin Flyway, un instrument care versionează modificările aduse structurii bazei de date prin fișiere SQL numerotate aplicate automat la pornirea aplicației, asigurând astfel că schema bazei de date este întotdeauna la versiunea corectă și că nicio migrare nu este aplicată de mai multe ori \cite{flyway}. Față de alte instrumente similare de versionare a schemei bazei de date, Flyway a fost preferat datorită simplității și integrării native cu Spring Boot, necesitând configurare minimă pentru a funcționa corect.
